\scanpage{145}
\mspage{3}
Cette catégorie s'interprète comme celle des triples $(x',\zeta,u)$ avec
\[
\left\{\begin{aligned}
 &x'\in\operatorname{Ob}C'\\
 &u:\unsure{g|f'}(x')=f(\alpha(x'))\longrightarrow s_0\\
 &\zeta:\deleted{\unsure{\alpha(x')}}x_0\longrightarrow\alpha(x')
\end{aligned}\right.
\]
de façon à rendre commutatif le diagramme
\diagram{\begin{tikzcd}[ampersand replacement=\&]
f(x_0)\arrow[rr,"u_0"]\arrow[dr,"f(\zeta)"']\&\&s_0\\
\&f(\alpha(x'))\arrow[ur,"u"]\&
\end{tikzcd}}

Je voudrais dégager des variantes de ces énoncés dans le cadre cohomologique, les foncteurs $\varprojlim$ étant remplacés par des $\unsure{\mathbb R|R}\Gamma$ (de préfaisceaux abéliens \inserted{p. ex.} -- mais on pense aussi aux ``cas non \unsure{commutatifs|commutatif}, définis par des dérivateurs \unsure{notamment}''.

% slanted
\margin{\unsure{(|?} \unsure{s|3}}
Si $I$ est une catégorie \uline{ordonnée} \uline{finie}, la condition ``filtrante à g.'' \deleted{implique} \inserted{équivaut à} l'existence d'un objet initial, ce qui implique que $I$ est contractile. \deleted{\unsure{Dans les énoncés} \illegible{2}} \inserted{Dans les énoncés précédents,} il se pourrait que en bonne philosophie soit de remplacer partout ``fbt filtrante à g'' par ``contractile. \deleted{\unsure{Donc}} Si p. ex. $C$ est la cat. finale, alors remplacer la condition que $C'$ soit filtr. à g, par \unsure{$W$-asphéricité}. C'est en tous cas en accord avec l'esprit de l'axiome \unsure{Dér 5} (cas où $J$ est la cat. finale) En effet, dans le cas où $C=e$, la condition \unsure{pour|que} $\alpha$ soit anodin pour les $\varprojlim$ équivaut visiblement à dire que le foncteur $\alpha^*:M\unsure{-|\to}\underline{\mathrm{Hom}}(C'{}^\circ,M)$, associant à tt $\unsure{F|E}\in M$ le préfaisceau constant de val. $E$
% archive sheet number: 145
